\begin{enumerate}
    \item La charge est soumise à deux forces~: son poids $\vec{P}$ et la force
          $\vec{T}$ (tension du câble). D'après le principe d'inertie~:
          \[
              \vec{P}+\vec{T}=\vec{0}
              \quad\text{ou}\quad
              \vec{T}=-\vec{P}
              \quad\text{ou}\quad
              T=P.
          \]

          Application numérique~:
          \[
              T=Mg=800\times10=\qty{8.0e3}{\N}.
          \]
    
    \item La puissance de la grue est égale à la puissance du travail de la
          force $\vec{T}$~:
    
          \[
              \mathcal{P}=\frac{W(\vec{T})}{\Delta t}=\vec{T}\cdot \vec{v}=Tv.
          \]
    
    
          Application numérique~:
    
          \[
              \mathcal{P}=8,0\times10^3\times2,0=\qty{1.6e4}{\W}=\qty{16}{\kW}.
          \]
\end{enumerate}

